
\section{Folio 3}

\begin{ejercicio}
    Sean $X,Y$ espacios de Banach, $T$ un operador lineal tal que $T: X \longrightarrow Y$, y $\{f_i\}_{i \in \mathcal{I}} \subseteq Y^*$ una familia que ``separa los puntos de $Y$'', es decir
    \begin{equation*}
        \forall y_1 \neq y_2, \ \exists j \in \mathcal{I} \text{ tal que } f_j(y_1) \neq f_j(y_2)
    \end{equation*}
    Demostrar que si $f_i \circ T$ es continuo $\forall i \in \mathcal{I}$, entonces $T$ es continuo. (\textit{Sugg.: utilizar el teorema del gráfico cerrado})\\ \\
    \noindent
    Vamos a usar el Teorema del Gráfico Cerrado, para ello entonces debemos demostrar que $T$ es cerrada, lo que equivale a demostrar que el gráfico $\Gamma_T$ es cerrado, es decir
    \begin{equation*}
        \forall (x_n) \subseteq X\ : \ x_n \longrightarrow x \text{ y } T(x_n) \longrightarrow y \implies y = T(x)
    \end{equation*}
    Supondremos las hipótesis, es decir, una sucesión cualquier $(x_n)$ que cumple
    \begin{equation*}
        x_n \longrightarrow x \text{ y } T(x_n) \longrightarrow y
    \end{equation*}
    entonces
    \begin{equation*}
        \lim_{n \to \infty} f_i(T(x_n)) \overset{(*)}{=} f_i\left(\lim_{n \to \infty} T(x_n)\right) = f_i(y)
    \end{equation*}
    donde en $(*)$ hemos usado que $f_i$ continua por pertenecer a $Y^*$. Mirando el Diagrama \ref{fig:diagrama-ejercicio1}, sabemos que $f_i \circ T$ es una aplicación de $X$ en $\mathbb{R}$ (y continua por hipótesis), y recordando que $x_n \longrightarrow x$, entonces
    $$\lim_{n \to \infty} (f_i \circ T)(x_n) = (f_i \circ T)(\lim_{n \to \infty} x_n) = (f_i\circ T)(x)=f_i(T(x))$$
    Uniendo ambos resultados tenemos que 
    \begin{equation*}
        f_i(T(x)) = f_i(y) \quad \forall i \in \mathcal{I}
    \end{equation*}
    Lo que queda ahora es pensar que aún siendo las imágenes iguales, ¿es posible que $y$ y $T(x)$ sean vectores distintos? Esto sabemos que no es posible pues en el caso de $y\neq T(x)$, al tener la propiedad de que la familia ``separa los puntos de $Y$'', entonces tendríamos 
    \begin{equation*}
        f_j(y) \neq f_j(T(x)) \quad \text{para algún } j \in \mathcal{I}
    \end{equation*}
    lo que es un absurdo con lo obtenido. Queda por tanto demostrado que
    \begin{equation*}
        y = T(x)
    \end{equation*}

\begin{figure}
    \centering
    % Definimos el diagrama con TikZ
    \begin{tikzpicture}[baseline=(X.base), every node/.style={on grid}, node distance=2cm]
        % Posicionamiento de los elementos en línea recta
        \node (X) at (0,0) {$X$};
        \node (Y) at (2,0) {$Y$};
        \node (R) at (4,0) {$\mathbb{R}$};
        
        % Flechas superiores rectas
        \draw[->] (X) -- node[above] {$T$} (Y);
        \draw[->] (Y) -- node[above] {$f_i$} (R);
        
        % Flecha inferior curva (semicírculo por abajo) desde X hasta R
        % "bend right=50" controla la curvatura hacia abajo
        \draw[->] (X) to[bend right=50] node[below=2pt] {$f_i \circ T$} (R);
    \end{tikzpicture}
    \caption{Diagrama de la aplicación $f_i \circ T$.}
    \label{fig:diagrama-ejercicio1}
\end{figure}

\end{ejercicio}


\begin{ejercicio}
    Sea $\{e_n\}_{n \in \mathbb{N}^+}$ una base ortonormal en un espacio de Hilbert separable $H$ y $T: H \longrightarrow H$ una aplicación lineal y continua. Demostrar que 
    \begin{equation*}
        T \text{ es isometría} \iff \{T(e_n)\}_{n \in \mathbb{N}^+} \text{ es un sistema ortonormal en } H
    \end{equation*}
    \begin{description}
        \item[$\Longrightarrow)$] Para probar que $\{T(e_n)\}_{n \in \mathbb{N}^+}$ es un sistema ortonormal, por definición, debemos demostrar que se cumplen dos condiciones para cualquier par de índices $n, m$:
        \begin{itemize}
            \item[$(i)$] $\|T(e_n)\| = 1$ para todo $n$ (norma de cada vector es 1).
            \item[$(ii)$] $( T(e_n), T(e_m) ) = 0$ para todo $n \neq m$ (vectores son ortogonales).
        \end{itemize}
        Sabiendo que $T$ es isometría veámos lo siguiente
        \begin{align*}
            \|x\|^2+2(x,y)+\|y\|^2 &= \|x+y\|^2 = \|T(x+y)\|^2 = (T(x+y), T(x+y)) =\\
            &= 2(T(x), T(y)) + \|T(x)\|^2 + \|T(y)\|^2 =\\
            &= 2(T(x), T(y)) + \|x\|^2 + \|y\|^2\implies (x,y) = (T(x), T(y)) \quad \forall x,y \in H
        \end{align*}
        Aplicando lo demostrado anteriormente tenemos que dados dos elementos cualesquiera de nuestra base ortonormal, $e_n$ y $e_m$, entonces 
        \begin{equation*}
            (T(e_n), T(e_m)) = (e_n, e_m) = 
            \begin{cases}
                1 & \text{si } n = m \\
                0 & \text{si } n \neq m
            \end{cases}
        \end{equation*}
        donde queda demostrado $(i)$ y $(ii)$, y por tanto $\{T(e_n)\}_{n \in \mathbb{N}^+}$ es un sistema ortonormal en $H$.
        
        \item[$\Longleftarrow)$] Aqui tenemos que demostrar que $T$ es una isometría, es decir, $\|T(x)\| = \|x\|$ para todo $x \in H$. \\
        \noindent
        Sea $x \in H$ un vector cualquiera del espacio de Hilbert. Como $\{e_n\}_{n \in \mathbb{N}^+}$ es una base ortonormal (que equivale a ser sistema ortonormal completo), podemos expandir el vector $x$ en su serie de Fourier abstracta
        
        $$x = \sum_{n=1}^{+\infty} \alpha_n e_n, \quad \text{donde } \alpha_n = ( x, e_n )$$
        y por la identidad de parseval se tiene
        \begin{equation}\label{eq:parseval}
            \|x\|^2 = \sum_{n=1}^{+\infty} |\alpha_n|^2
        \end{equation}
        Introduciendo la igualdad de cualquier $x\in \mathcal{H}$ en la aplicación de $T$ tenemos
        \begin{equation*}
            T(x) = T\left( \sum_{n=1}^{+\infty} \alpha_n e_n \right) \overset{(*)}{=} \sum_{n=1}^{+\infty} \alpha_n T(e_n)\implies \|T(x)\|^2 = \left( \sum_{n=1}^{+\infty} \alpha_n T(e_n), \; \sum_{m=1}^{+\infty} \alpha_m T(e_m) \right)
        \end{equation*}
        donde en $(*)$ hemos aplicado linealidad y continuidad (esta última para poder meter el operador dentro del límite de la sumatoria). Siguiendo con el mismo desarrollo tenemos
        \begin{equation*}
            \|T(x)\|^2 = \sum_{n=1}^{+\infty} \sum_{m=1}^{+\infty} \alpha_n \overline{\alpha_m} ( T(e_n), T(e_m) )= \sum_{n=1}^{+\infty} |\alpha_n|^2 \cdot 1 = \sum_{n=1}^{+\infty} |\alpha_n|^2=\|x\|^2
        \end{equation*}
        donde hemos demostrado que $T$ es una isometría, y por tanto queda demostrado el ejercicio.
    \end{description}
\end{ejercicio}

\begin{ejercicio} % //TODO: hacer
    Sea $M$ un subespacio cerrado de uno espacio de Hilbert $H$ y $T \in BL(H)$. Demostrar que:
    \begin{itemize}
        \item[i)] si $T$ es autoadjunto y $M$ es invariante por $T$ (es decir, $T(M) \subseteq M$), entonces también $M^\perp$ es invariante por $T$ (es decir, $T(M^\perp) \subseteq M^\perp$).
        
        \item[ii)] si $M$ y $M^\perp$ son invariantes por $T$, entonces lo son también por $T^*$.
    \end{itemize}
\end{ejercicio}


\begin{ejercicio}
    Sea $T: C^0([0,\pi]) \longrightarrow C^0([0,\pi])$ el siguiente operador lineal:
    \[ T(u)(t) = \sin(t) u(t), \qquad t \in [0,\pi] \]
    Encontrar la norma de $T$, el espectro de $T$, los posibles valores propios y los correspondientes espacios propios. ¿Es el operador $T$ compacto? \\
    \noindent
    Recordemos que aqui trabajamos con la norma del supremo dada así
    $$\|u\|_\infty =\sup_{t \in [0,\pi]} |u(t)|=\max_{t \in [0,\pi]} |u(t)|$$
    donde trabajamos con el máximo (en lugar de supremo), porque al ser funcioines contiuas, el teorema de Weierstrass nos asegura que se alcanza el máximo. Sabemos que $C^0([0,\pi])$ es el espacio de las funciones continuas en $[0,\pi]$ que es un espacio de Banach. \\

    \noindent
    Calcularemos ahora la \underline{norma} del supremo, para ello primero buscaremos la acotación susperior, veámoslo
    $$\|T(u)\|_\infty = \max_{t \in [0,\pi]} |T(u)(t)| = \max_{t \in [0,\pi]} |\sin(t) u(t)|$$
    Sabemos que para cualquier $t \in [0,\pi]$, el seno está acotado por $1$ ($|\sin(t)| \le 1$). Por lo tanto, podemos separarlo hacia arriba:$$\|T(u)\|_\infty \le \max_{t \in [0,\pi]} (1 \cdot |u(t)|) = \max_{t \in [0,\pi]} |u(t)| = \|u\|_\infty \quad \forall u\in C^0([0,\pi])$$
    Pasando la norma de $u$ dividiendo, el supremo de la fracción no puede superar a $1$, por lo que
    $$\|T\| = \sup_{u \neq 0} \frac{\|T(u)\|_\infty}{\|u\|_\infty} \leq 1$$
    Ahora falta tomar un $u\in C^0([0,\pi])$ tal que encontremos que esa fracción es $1$, para ello tomaremos $u(t)=1 \ \forall t\in [0,\pi]$, y entonces tenemos
    \begin{equation*}
        \|T(u)\|_\infty = \max_{t \in [0,\pi]} |\sin(t) \cdot 1| = \max_{t \in [0,\pi]} |\sin(t)| = 1
    \end{equation*}
    donde el máximo se alcanza en $t = \frac{\pi}{2}$, y sabiendo que $\|1\|_\infty = 1$, tenemos que
    $$\|T\| = \sup_{u \neq 0} \frac{\|T(u)\|_\infty}{\|u\|_\infty} \geq \frac{\|T(1)\|_\infty}{\|1\|_\infty} = 1\implies \|T\|=1$$
    Estudiaremos ahora los \underline{valores propios} y los \underline{subespacios propios asociados}. Planteamos la ecuación de autovalores $T(u) = \lambda u$ para una función no nula $u \neq 0$. Esto significa que en cada punto $t \in [0,\pi]$ debe cumplirse
    $$\sin(t) u(t) = \lambda u(t) \implies \big(\sin(t) - \lambda\big) u(t) = 0 \quad \forall t \in [0,\pi]$$
    Para que $\lambda$ sea un valor propio, debe existir una función continua $u(t)$ que no sea idénticamente cero y cumpla esa igualdad en todo el intervalo. Analicemos qué ocurre:
    \begin{itemize}
        \item Si fijamos un valor de $\lambda$, el paréntesis $\big(\sin(t) - \lambda\big)$ solo puede valer cero, a lo sumo, en uno o dos puntos aislados del intervalo $[0,\pi]$ (por ejemplo, si $\lambda = 1$, solo vale cero en $t = \pi/2$).
        \item En todos los demás puntos de la línea donde $\sin(t) \neq \lambda$, el paréntesis no es cero, lo que obliga matemáticamente a que $u(t) = 0$.
        \item Como la función $u(t)$ debe ser obligatoriamente continua, si vale cero en casi todo el intervalo, por continuidad se ve arrastrada a valer cero también en esos puntos aislados.
    \end{itemize}
    Por tanto, la única función continua que satisface la ecuación para cualquier $\lambda$ es la función nula $u(t) = 0$. Como no existe ninguna función testigo no nula, concluimos que el operador no tiene ningún valor propio, es decir
    \begin{equation*}
        e(T)=\emptyset
    \end{equation*}
    Calculemos ahora el \underline{espectro} de $T$, donde ya sabemos que

    \begin{equation*}
        \sigma(T)\subseteq \{\lambda \in \C\mid |\lambda| \leq \|T\|\} = \{\lambda \in \C\mid |\lambda| \leq 1\}
    \end{equation*}
    Por definición, un número $\lambda \in \mathbb{C}$ pertenece al espectro $\sigma(T)$ si el operador $(T - \lambda I)$ no es biyectivo (es decir, o falla la inyectividad o falla la sobreyectividad). Para comprobar esto, planteamos la ecuación fundamental del espectro: intentamos despejar una función incógnita $u(t)$ a partir de una función dada $v(t)$ (ambas continuas en $[0,\pi]$)
    $$(T - \lambda I)u(t) = v(t) \implies \sin(t)u(t) - \lambda u(t) = v(t)$$
    Si queremos despejar nuestra función solución $u(t)$, el álgebra nos obliga a pasar dividiendo el paréntesis
    $$u(t) = \frac{v(t)}{\sin(t) - \lambda}$$
    A partir de esta fracción, el análisis se divide en dos mundos completamente opuestos dependiendo de si el denominador se puede hacer cero o no.
    \begin{itemize}
        \item \tbf{Caso $\lambda\notin [0,1]$}: en este caso, el denominador $\sin(t) - \lambda$ no se anula para ningún \mbox{$t \in [0,\pi]$}. Esto significa que la función solución $u(t)$ es continua en todo el intervalo, ya que el denominador no se acerca a cero y el numerador $v(t)$ es continua (se aplica la sobreyectividad), y como es única punto a punto por que tenemos su definición (aplica inyectividad), tenemos la función buscada. Por lo tanto, si $\lambda \notin [0,1]$, entonces $\lambda \in \rho(T)$ (resolvente).
        \item \tbf{Caso $\lambda \in [0,1]$}: en este caso, el denominador $\sin(t) - \lambda$ se anula en al menos un punto del intervalo $[0,\pi]$, es decir, donde el seno valga exactamente $\lambda$
        $$\sin(t_0) = \lambda \implies \sin(t_0) - \lambda = 0$$
        En este sentido, si tomamos una función $v\in C^0([0,\pi])$ donde $v(\lambda)\neq 0$, el denominador se anula en $t_0$ y la función solución $u(t)$ se dispara al infinito, lo que rompe la continuidad de $u(t)$ (falla la sobreyectividad). Por lo tanto, si $\lambda \in [0,1]$, entonces $\lambda \in \sigma(T)$.
    \end{itemize}
    Concluimos por tanto que
    \begin{equation*}
        \sigma(T) = [0,1]
    \end{equation*}
    Para demostrar si $T$ es compacto o no, es como suponer que sí lo es, entonces tendríamos la siguiente propiedad de los operadores compactos, lineales y continuos

    \begin{equation*}
        \sigma(T) \setminus \{0\}= e(T) \setminus \{0\}
    \end{equation*}
    como esto es una contradicción, y sabemos que $T\in BL(C^0([0,\pi]))$, entonces $T$ no es compacto. Recordemos que la linealidad es trivialmente comprobable, y la continuidad viene dada por saber cual es la norma del operador y saber por tanto que está acotado.
\end{ejercicio}

\begin{ejercicio}
    Sea $T$ el siguiente operador lineal:
    \[ T(x) = (\alpha_1 x_1, \alpha_2 x_2, \dots, \alpha_k x_k, \dots), \qquad \forall x = (x_1, \dots, x_k, \dots) \in \ell_2 \]
    donde $\alpha_k = \dfrac{k}{k+1}, \quad \forall k \in \mathbb{N}^+$. Demostrar que $T(x) \in \ell_2, \ \forall x \in \ell_2$ y que $T$ es continuo. \\
    Encontrar después todos los valores propios de $T$ y los espacios propios correspondientes. ¿Es el operador $T$ compacto? \\

    \noindent
    Para demostrar que $T(x) \in \ell_2$ para cualquier $x \in \ell_2$, debemos verificar que la serie de los cuadrados de las componentes de $T(x)$ converge, es decir, que
    \begin{equation*}
        \sum_{k=1}^{\infty} |\alpha_k x_k|^2 < \infty
    \end{equation*}
    para ello veámos el valor de cada $\alpha_k$
    \begin{equation}\label{eq:alpha}
        |\alpha_k| = \frac{k}{k+1} < 1 \quad \forall k \in \mathbb{N}^+
    \end{equation}
    Veámos ahora el desarrollo que buscamos

    \begin{equation*}
        \|T(x)\|^2 = \sum_{k=1}^{+\infty} |\alpha_k x_k|^2 = \sum_{k=1}^{+\infty} |\alpha_k|^2 |x_k|^2\overset{(*)}{\leq} \sum_{k=1}^{+\infty} |x_k|^2 = \|x\|^2 < \infty
    \end{equation*}
    donde en $(*)$ hemos usado lo obtenido en (\ref{eq:alpha}), y sabiendo $\|x\|^2<\infty$ porque $x\in \ell_2$ tenemos lo buscado, es decir, $T(x)\in \ell_2$. \\

    \noindent
    Para demostrar que $T$ es continuo, basta con demostrar que $T$ es acotado, es decir, que existe una constante $C > 0$ tal que $\|T(x)\| \leq C \|x\|$ para todo $x \in \ell_2$. En la Parte 1 acabamos de obtener directamente que
    $$\|T(x)\|^2 \le \|x\|^2 \implies \|T(x)\| \le 1 \cdot \|x\| \quad \forall x \in \ell_2$$
    donde queda demostrada la \underline{continuidad} de $T$. \\

    \noindent
    Nos centraremos ahora en buscar los \underline{valores propios} y los \underline{subespacios propios asociados}. Por definición, un número complejo $\lambda \in \mathbb{C}$ es un valor propio de $T$ si existe un vector no nulo $x \neq 0$ en $\ell_2$ tal que
    $$T(x) = \lambda x$$
    Planteamos la ecuación componente a componente
    $$(\alpha_1 x_1, \alpha_2 x_2, \dots, \alpha_k x_k, \dots) = (\lambda x_1, \lambda x_2, \dots, \lambda x_k, \dots) \implies (\alpha_k - \lambda) x_k = 0 \quad \forall k \in \mathbb{N}^+$$
    Para cada componente $k$, para que esta ecuación se cumpla, tenemos dos opciones: o bien $x_k = 0$, o bien $\lambda = \alpha_k$. Como buscamos un vector propio $x \neq 0$, al menos una componente tiene que ser distinta de cero (digamos, la componente número $j$, $x_j \neq 0$). Si $x_j \neq 0$, la ecuación nos obliga a que
    $$\lambda = \alpha_j = \frac{j}{j+1}$$
    Por lo tanto, los valores propios de $T$ son exactamente los elementos de la sucesión $\{\alpha_k\}$
    $$\lambda_k = \frac{k}{k+1} \quad \forall k \in \mathbb{N}^+\implies e(T)=\left\{\dfrac{k}{k+1} : k \in \mathbb{N}^+\right\}$$
    Si fijamos un valor propio $\lambda_k$, sabemos que el vector $x\in \ell_2$ en la componente $k$ puede tomar cualquier valor, pero para el resto de componentes $j\neq k$ debemos tener la igualdad
    \begin{equation*}
        (\alpha_j - \lambda_k) x_j = 0 \implies x_j = 0 \quad \forall j \neq k
    \end{equation*}
    de donde deducimos que para cada componente $k$, el espacio propio asociado a $\lambda_k$ es
    $$V_k = \ker(T - \lambda_k I) = \text{span}\{e_k\}= \{ (0, \dots, 0, c, 0, \dots) : c \in \mathbb{C} \}$$
    Comprobemos ahora si $T$ es \underline{compacto} o no. Para ello, vamos a demostrar que no es commpacto usando la base ortonormal canónica de $\ell_2$, $\{e_n\}_{n \in \mathbb{N}^+}$, la cual sabemos que está acotada porque $\|e_n\| = 1$ para todo $n$. \\

    \noindent
    Si aplicamos el operador $T$ a los elementos de esta base, dado que $e_n = (0, \dots, 1, 0, \dots)$ tiene un $1$ en la posición $n$, el operador simplemente lo multiplica por su correspondiente $\alpha_n$
    $$T(e_n) = \alpha_n e_n = \frac{n}{n+1} e_n$$
    por lo que debemos demostrar que de esta sucesión de imágenes $\{T(e_n)\}$ es imposible extraer ninguna subsucesión que converja fuertemente (en norma). Calculemos la distancia al cuadrado entre dos términos distintos cualesquiera de la sucesión de imágenes, $T(e_n)$ y $T(e_m)$ (con $n \neq m$)
    \begin{align*}
        \|T(e_n) - T(e_m)\|^2 &= \|(\alpha_n - \alpha_m) e_n + \alpha_m (e_n - e_m)\|^2= \\
        &=(\alpha_n e_n - \alpha_m e_m, \; \alpha_n e_n - \alpha_m e_m )=|\alpha_n|^2\|e_n\|^2 + |\alpha_m|^2 \|e_m\|^2=\\
        &= \left(\frac{n}{n+1}\right)^2 + \left(\frac{m}{m+1}\right)^2 
    \end{align*}
    donde si pasamos a cuando $k,m\to \infty$ tenemos que
    \begin{equation*}
        \lim_{n,m \to \infty} \|T(e_n) - T(e_m)\|^2 = \lim_{n,m \to \infty} \left(\frac{n}{n+1}\right)^2 + \left(\frac{m}{m+1}\right)^2 = 2\implies \lim_{n,m \to \infty} \|T(e_n) - T(e_m)\| = \sqrt{2} \neq 0
    \end{equation*}
    lo que demuestra que los términos de la sucesión de imágenes, a medida que avanzamos en el infinito, se estabilizan a una distancia fija de $\sqrt{2}$ entre sí, y por tanto no es posible extraer ninguna subsucesión que converja fuertemente, lo que implica que $T$ \underline{no es compacto}.
\end{ejercicio}


\begin{ejercicio}
    Sea $T: \ell_2 \longrightarrow \ell_2$ definido por:
    \[ T(x) = \left( x_1, x_2, 0, \frac{x_3}{2^3}, \frac{x_4}{2^4}, \dots, \frac{x_n}{2^n}, \dots \right) \]
    $\forall x = (x_1, \dots, x_n, \dots) \in \ell_2$.
    \begin{itemize}
        \item[i)] Demostrar que $T$ es compacto
        \item[ii)] Determinar el espectro de $T$, los posibles valores propios y los correspondientes espacios propios.
    \end{itemize}
    Antes de resolver el ejercicio, es importante notar que $T$ es lineal y continuo, donde la linealidad es trivial y la continuidad se puede demostrar fácilmente usando la acotación de la siguiente manera. \\

    \noindent
    Recordemos que ser acotado significa que
    \begin{equation*}
        \exists C > 0 : \|T(x)\|_2 \leq C \|x\|_2 \quad \forall x \in \ell_2
    \end{equation*}
    Vamos a calcular la norma de $T(x)$ al cuadrado
    $$\|T(x)\|_2^2 = |x_1|^2 + |x_2|^2 + 0^2 + \sum_{n=3}^{\infty} \left| \frac{x_n}{2^{n}} \right|^2$$
    Como todos los denominadores $2^{n}$ son mayores o iguales que 1, es evidente que cada fracción $\dfrac{1}{(2^{n})^2} \le 1$. Por lo tanto, podemos quitar los denominadores y la desigualdad se mantiene hacia arriba:$$\|T(x)\|_2^2 \le |x_1|^2 + |x_2|^2 + \sum_{n=3}^{\infty} |x_n|^2 = \sum_{n=1}^{\infty} |x_n|^2 = \|x\|_2^2$$Tomando la raíz cuadrada en ambos lados obtenemos:$$\|T(x)\|_2 \le 1 \cdot \|x\|_2 \implies \|T\| \le 1$$Como la norma del operador está acotada (es finita), $T$ es un operador continuo.
    \begin{description}
        \item[$i)$] Vamos a demostrar la compacidad de $T$ aproximandolo por operadores de rango finito. Para cada $k \geq 3$, definimos el operador $T_k$ que copia a $T$ hasta la componente $k$ y a partir de ahí pone ceros
        $$T_k(x) = \left( x_1, x_2, 0, \frac{x_3}{2^3}, \dots, \frac{x_{k-1}}{2^{k-1}}, 0, 0, \dots \right)$$
        Como la imagen de $T_k$ tiene como máximo dimensión $k$ (dimensión finita), $T_k$ es un operador de rango finito y, por tanto, es automáticamente compacto. Restamos ambos operadores
        $$(T - T_k)(x) = \left( 0, \dots, 0, \frac{x_k}{2^{k}}, \frac{x_{k+1}}{2^{k+1}}, \dots \right)$$
        Evaluamos su norma en $\ell_2$ al cuadrado
        $$\|(T - T_k)(x)\|_2^2 = \sum_{n=k}^{\infty} \left| \frac{x_n}{2^{n}} \right|^2 = \sum_{n=k}^{\infty} \frac{1}{(2^{n})^2} |x_n|^2 \leq \frac{1}{(2^{k})^2} \sum_{n=k}^{\infty} |x_n|^2$$
        Como la última suma es menor o igual que la norma total $\|x\|_2^2$, acotamos
        $$\|(T - T_k)(x)\|_2^2 \leq \frac{1}{4^{k}} \|x\|_2^2 \implies \|(T - T_k)(x)\|_2 \leq \frac{1}{2^{k}} \|x\|_2$$
        Por la definición de la norma de un operador
        $$\|T - T_k\| = \sup_{x \neq 0} \frac{\|(T - T_k)(x)\|_2}{\|x\|_2} \leq \frac{1}{2^{k}}$$
        Si tomamos el límite cuando $k \to \infty$
        $$\lim_{k \to \infty} \|T - T_k\| = 0$$
        Como $T$ es el límite en norma de una sucesión de operadores compactos ($T_k$) y el espacio de operadores compactos es cerrado, $T$ es un operador compacto, como se quería demostrar.

        \item[$ii)$] Primero intentaremos obtener los posibles \underline{valores propios} de $T$, y sus \underline{subespacios propios asociados}. Para ello planteamos la ecuación general $(T - \lambda I)x = 0$ componente a componente
        \begin{equation*}
            \begin{cases}
                x_1 - \lambda x_1 = 0 \implies (1 - \lambda)x_1 = 0 \\
                x_2 - \lambda x_2 = 0 \implies \left(1 - \lambda\right)x_2 = 0 \\
                0 - \lambda x_3 = 0 \implies -\lambda x_3 = 0 \\
                \dfrac{x_{n-1}}{2^{n-1}} - \lambda x_n = 0 \quad \text{para } n \geq 4
            \end{cases}
        \end{equation*}
        Recordemos una cosa de vital importancia, lo que realmente estamos buscando son los valores de $\lambda$ para los cuales existe un vector $x \neq 0$ que cumple esta ecuación. Es decir, estamos buscando los valores de $\lambda$ para los cuales el núcleo $\ker(T - \lambda I)$ tiene más elementos que el cero. \\

        \noindent
        Dicho esto, veámos los distintos casos para $\lambda$:
        \begin{itemize}
            \item \tbf{Caso $\lambda = 1$}: de las ecuaciones (1) y (2) obtenemos $0 \cdot x_1 = 0$ y $0 \cdot x_2 = 0$. Esto significa que $x_1$ y $x_2$ pueden ser cualquier número. Las demás componentes se obligan a ser 0. Por tanto, $\lambda = 1$ es un valor propio. El espacio propio asociado es:
            $$V_1 = \ker(T - I) = \text{span}\{(1,0,0,\dots), (0,1,0,\dots)\} = \text{span}\{e_1, e_2\}$$
            \item \tbf{Caso $\lambda = 0$}: de la ecuación (3) obtenemos $0 \cdot x_3 = 0$. Esto significa que $x_3$ puede ser cualquier número. Sustituyendo en las demás ecuaciones con $\lambda = 0$, tenemos lo siguiente:
            \begin{itemize}
                \item De (1): $x_1 = 0$.
                \item De (2): $x_2 = 0$.
                \item De (4): $\dfrac{x_3}{2^3} = 0 \implies x_3 = 0$, por lo que en esta ecuación, se nos obliga a que $x_3$ también sea 0.
            \end{itemize}
            Si continuamos despejando, veremos que todas las componentes se hacen cero, por lo que  para $\lambda = 0$, el único vector que cumple la ecuación es el $x = 0$, y por tanto $0$ NO es un valor propio.
            \item \tbf{Caso $\lambda \neq 1$ y $\lambda \neq 0$}: de (1) y (2) se deduce que $x_1 = 0$ y $x_2 = 0$. De (3) se deduce que $x_3 = 0$. Usando la fórmula recursiva (4), como todas las anteriores son cero, todas las componentes siguientes se anulan. Solo existe la solución trivial $x=0$.
        \end{itemize}
        El único valor propio es $\sigma_p(T) = \{1\}$, y su espacio propio asociado es 
        $$V_1 = \text{span}\{e_1, e_2\}$$
        Ahora nos centraremos en obtener el \tbf{espectro} de $T$, donde de primeras ya sabemos que 
        \begin{equation*}
            e(T) = \{1\} \subseteq \sigma(T)
        \end{equation*}
        y además por ser $\ell_2$ dimensión infinita y $T$ compacto sabemos que $0\in \sigma(T)$. \\

        \noindent
        Tomemos ahora $\lambda\neq0,1$ y vamos a intentar despejar de forma única el vector $x$ en función de $y$, recordando primero que la ecuación general es $(T - \lambda I)x = y$, es decir, que
        \begin{equation*}
            \left( \left(1 - \lambda\right)x_1, \left(1 - \lambda\right)x_2, -\lambda x_3, \left(\frac{x_{3}}{2^{3}} - \lambda x_4\right), \dots, \left(\frac{x_{n-1}}{2^{n-1}} - \lambda x_n\right), \dots \right) = (y_1, y_2, y_3, y_4, \dots)
        \end{equation*}
        Despejando cada componente $x_n$ de forma única, obtenemos:
        \begin{itemize}
            \item $x_1 = \dfrac{y_1}{1 - \lambda}$
            \item $x_2 = \dfrac{y_2}{1 - \lambda}$
            \item $x_3 = \dfrac{y_3}{-\lambda}$
            \item $x_n = \dfrac{1}{\lambda} \left( \dfrac{x_{n-1}}{2^{n-1}} - y_n \right) \quad \text{para } n \geq 4$
        \end{itemize}
        como $\lambda \neq 1$ y $\lambda \neq 0$, los denominadores nunca se anulan (el despeje es único, luego es inyectivo). Además, debido a que dividimos sucesivamente por $2^{n-1}$, los términos $x_n$ decrecen rapidísimo geométricamente. Eso garantiza que para cualquier $y \in \ell_2$, la serie de $x$ va a converger, quedándose dentro de $\ell_2$ (es sobreyectivo). Al ser inyectivo y sobreyectivo, el operador es biyectivo, por lo que todos estos valores de $\lambda$ están en el resolvente $\rho(T)$, lo que significa que el espetro es
        \begin{equation*}
            \sigma(T) = \{0,1\}
        \end{equation*}
    \end{description}
\end{ejercicio}

